% ── §1.1: Αντικείμενο και Κίνητρο ───────────────────────────

\section{Αντικείμενο και Κίνητρο}
\label{sec:motivation}

Η γήρανση του πληθυσμού σε συνδυασμό με την αυξανόμενη επικράτηση χρόνιων
παθήσεων δημιουργεί έντονες πιέσεις στα συστήματα υγειονομικής περίθαλψης
παγκοσμίως. Τα περιβάλλοντα Υποβοηθούμενης Διαβίωσης (Ambient Assisted
Living, AAL) αναδύονται ως μία από τις πλέον υποσχόμενες τεχνολογικές
λύσεις, παρέχοντας στους ηλικιωμένους και στους ανθρώπους με ειδικές
ανάγκες τα μέσα για ανεξάρτητη διαβίωση με ταυτόχρονη συνεχή
παρακολούθηση της κατάστασης υγείας τους \cite{Liu2021}.

Κεντρικό στοιχείο των συστημάτων AAL αποτελεί η Αναγνώριση Δραστηριότητας
Ανθρώπου (Human Activity Recognition, HAR), η οποία επιδιώκει την αυτόματη
αναγνώριση καθημερινών κινητικών δραστηριοτήτων — όπως το περπάτημα, η
ανάβαση σκαλιών, το κάθισμα και η ανάπαυση — μέσω δεδομένων αισθητήρων
έξυπνων συσκευών. Η πρώιμη ανίχνευση ανωμαλιών στη δραστηριότητα, όπως
αιφνίδιες πτώσεις ή παρατεταμένη ακινησία, επιτρέπει την έγκαιρη
παρέμβαση κλινικού προσωπικού, βελτιώνοντας σημαντικά την ποιότητα ζωής
και μειώνοντας τον κίνδυνο σοβαρών επιπλοκών \cite{Liu2021}.

Παρά την αποτελεσματικότητά τους, οι παραδοσιακές προσεγγίσεις μηχανικής
μάθησης για HAR βασίζονται στο κεντρικό παράδειγμα (centralized paradigm),
σύμφωνα με το οποίο τα ακατέργαστα δεδομένα αισθητήρων μεταφέρονται σε
κεντρικούς εξυπηρετητές για εκπαίδευση μοντέλων. Αυτή η προσέγγιση
συνεπάγεται σοβαρούς κινδύνους για την ιδιωτικότητα: τα δεδομένα κίνησης
και δραστηριότητας φέρουν ευαίσθητες πληροφορίες για τη φυσική κατάσταση,
τις συνήθειες και την υγεία του ατόμου \cite{Shokri2015}. Ερευνητικές
εργασίες έχουν αποδείξει ότι μέσω αναλυτικών επιθέσεων ανάκτησης δεδομένων
(inference attacks) είναι δυνατή η εξαγωγή ευαίσθητων χαρακτηριστικών
ακόμα και από τις ενημερώσεις παραμέτρων μοντέλου, εάν αυτές δεν
προστατευθούν κατάλληλα \cite{Shokri2015}.

Η Ομοσπονδιακή Μάθηση (Federated Learning, FL) \cite{McMahan2017} προσφέρει
μία θεμελιωδώς διαφορετική προσέγγιση: το μοντέλο εκπαιδεύεται τοπικά σε
κάθε συσκευή χρήστη και μόνο οι ενημερώσεις παραμέτρων (gradients ή βάρη)
αποστέλλονται στον κεντρικό εξυπηρετητή για συνάθροιση. Τα πρωτογενή
δεδομένα δεν εγκαταλείπουν ποτέ τη συσκευή, διασφαλίζοντας έτσι την
ιδιωτικότητα εκ σχεδιασμού (privacy by design). Ωστόσο, η πρακτική
εφαρμογή του FL σε AAL περιβάλλοντα αντιμετωπίζει δύο θεμελιώδεις
προκλήσεις που παραμένουν σε μεγάλο βαθμό ανεπίλυτες στη βιβλιογραφία.

Η \textbf{πρώτη πρόκληση} αφορά στην \textit{έλλειψη κινήτρων συμμετοχής}.
Σε πραγματικά συστήματα, οι χρήστες είναι ορθολογικοί πράκτορες που
αξιολογούν το κόστος συμμετοχής (κατανάλωση μπαταρίας, επεξεργαστική
ισχύς, εύρος ζώνης) έναντι του οφέλους που λαμβάνουν. Χωρίς κατάλληλους
μηχανισμούς κινήτρων, η συμμετοχή παραμένει χαμηλή ή συγκεντρώνεται σε
clients με πλεονάζοντες πόρους, υπονομεύοντας τη γενίκευση του μοντέλου
\cite{Kang2019}.

Η \textbf{δεύτερη πρόκληση} αφορά στην \textit{αδυναμία δίκαιης αξιολόγησης
συνεισφοράς}. Ακόμα και όταν υπάρχουν ανταμοιβές, ο ισοκατανεμημένος
υπολογισμός τους αγνοεί το γεγονός ότι διαφορετικοί clients συνεισφέρουν
άνισα στη βελτίωση του καθολικού μοντέλου. Clients με σπάνιες κλάσεις
δραστηριότητας ή με δεδομένα υψηλής ποιότητας θα πρέπει να ανταμείβονται
αναλόγως, ενώ clients με αδύναμη ή πλεονάζουσα συνεισφορά να λαμβάνουν
μικρότερη ανταμοιβή \cite{Wang2019}.

Η παρούσα εργασία απευθύνεται ακριβώς σε αυτό το κενό, προτείνοντας ένα
ολοκληρωμένο FL σύστημα για HAR σε AAL περιβάλλοντα που ενσωματώνει
ταυτόχρονα μηχανισμό κινήτρων και δίκαιη αξιολόγηση συνεισφοράς.
